\ifdefined\gkver\else\def\gkver{student}\fi
\documentclass[\gkver]{gaokaozhenti}
\begin{document}

\chapter{2004年江苏春季理综}
%% paperType: 理综物理部分

\begin{enumerate}
\item
%% number: 17
%% paperName: 2004年江苏春季理综第 17 题 3 分
%% typeId: 1
%% type: 单选题
%% chapter: 运动和力
%% point: 牛顿第三定律
%% method: 无
%% score: 3
%% degree: 850
%% duplicateId: 0
%% body:
关于两个物体间的作用力和反作用力，下列说法中正确的是\xzanswer{A}
\fourchoices[answer=A]
{作用力和反作用力一定同时产生，同时消失}
{作用力和反作用力可以不同时产生}
{作用力和反作用力可以是不同性质的力}
{作用力和反作用力的效果会相互抵消}
%% answer: A
%% memo:
\memoanswer{
本题原卷及材料中未提供详解，待补充。
}

\item
%% number: 18
%% paperName: 2004年江苏春季理综第 18 题 3 分
%% typeId: 1
%% type: 单选题
%% chapter: 气体性质
%% point: 物体的内能
%% method: 无
%% score: 3
%% degree: 850
%% duplicateId: 0
%% body:
关于分子的热运动，下列说法中正确的是\xzanswer{C}
\fourchoices[answer=C]
{当温度升高时，物体内每一个分子热运动的速率一定都增大}
{当温度降低时，物体内每一个分子热运动的速率一定都减小}
{当温度升高时，物体内分子热运动的平均动能必定增大}
{当温度降低时，物体内分子热运动的平均动能也可能增大}
%% answer: C
%% memo:
\memoanswer{
本题原卷及材料中未提供详解，待补充。
}

\item
%% number: 19
%% paperName: 2004年江苏春季理综第 19 题 3 分
%% typeId: 1
%% type: 单选题
%% chapter: 电场
%% point: 库仑定律
%% method: 无
%% score: 3
%% degree: 850
%% duplicateId: 0
%% body:
两个固定不动的点电荷，它们之间的静电力大小为 $F$，现使它们的电荷量都变为原来的一半，则它们之间的静电力大小变为\xzanswer{B}
\fourchoices[answer=B]
{$4F$}
{$F/4$}
{$2F$}
{$F/2$}
%% answer: B
%% memo:
\memoanswer{
本题原卷及材料中未提供详解，待补充。
}

\item
%% number: 20
%% paperName: 2004年江苏春季理综第 20 题 3 分
%% typeId: 1
%% type: 单选题
%% chapter: 光学
%% point: 光的电磁说
%% method: 无
%% score: 3
%% degree: 850
%% duplicateId: 0
%% body:
下列关于紫外线的几种说法中，正确的是\xzanswer{C}
\fourchoices[answer=C]
{紫外线是一种紫色的可见光}
{紫外线的频率比红外线的频率低}
{紫外线可使钞票上的荧光物质发光}
{利用紫外线可以进行电视机等电器的遥控}
%% answer: C
%% memo:
\memoanswer{
本题原卷及材料中未提供详解，待补充。
}

\item
%% number: 21
%% paperName: 2004年江苏春季理综第 21 题 3 分
%% typeId: 1
%% type: 单选题
%% chapter: 电场
%% point: 电场线
%% method: 无
%% score: 3
%% degree: 850
%% duplicateId: 0
%% body:
电场中某区域的电场线分布如图\ref{2004江苏春季理综21}所示，A、B 为电场中的两点。若以 $E_{A}$、$E_{B}$ 分别表示 A、B 两点的场强，$\varphi_{A}$、$\varphi_{B}$ 分别表示 A、B 两点的电势，则以下结论中正确的是\xzanswer{D}
\onepicture[num=21, numstyle=arabic, labelprefix={图}, label=2004江苏春季理综21, w=4.4cm]{figs/21.png}
\fourchoices[answer=D]
{$E_{A}>E_{B}$，$\varphi_{A}>\varphi_{B}$}
{$E_{A}<E_{B}$，$\varphi_{A}<\varphi_{B}$}
{$E_{A}>E_{B}$，$\varphi_{A}<\varphi_{B}$}
{$E_{A}<E_{B}$，$\varphi_{A}>\varphi_{B}$}
%% answer: D
%% memo:
\memoanswer{
本题原卷及材料中未提供详解，待补充。
}

\item
%% number: 22
%% paperName: 2004年江苏春季理综第 22 题 3 分
%% typeId: 1
%% type: 单选题
%% chapter: 磁场
%% point: 磁场的叠加
%% method: 无
%% score: 3
%% degree: 850
%% duplicateId: 0
%% body:
如图\ref{2004江苏春季理综22}所示，两根无限长的平行导线水平放置，两导线中均通以向右的、大小相等的恒定电流 $I$，图中的 A 点与两导线共面，且到两导线的距离相等。则这两根通电导线在该点产生的磁场的磁感应强度的合矢量\xzanswer{C}
\onepicture[num=22, numstyle=arabic, labelprefix={图}, label=2004江苏春季理综22, w=6.0cm]{figs/22.png}
\fourchoices[answer=C]
{方向水平向右}
{方向水平向左}
{大小一定为零}
{大小一定不为零}
%% answer: C
%% memo:
\memoanswer{
本题原卷及材料中未提供详解，待补充。
}

\item
%% number: 23
%% paperName: 2004年江苏春季理综第 23 题 3 分
%% typeId: 1
%% type: 单选题
%% chapter: 光学
%% point: 光的电磁说
%% method: 无
%% score: 3
%% degree: 850
%% duplicateId: 0
%% body:
关于机械波和电磁波，下列说法中错误的是\xzanswer{A}
\fourchoices[answer=A]
{机械波和电磁波都能在真空中传播}
{机械波和电磁波都可以传递能量}
{波长、频率和波速间的关系，即 $v=\lambda f$，对机械波和电磁波都适用}
{机械波和电磁波都能发生衍射和干涉现象}
%% answer: A
%% memo:
\memoanswer{
本题原卷及材料中未提供详解，待补充。
}

\item
%% number: 24
%% paperName: 2004年江苏春季理综第 24 题 3 分
%% typeId: 1
%% type: 单选题
%% chapter: 原子和原子核
%% point: 质能方程
%% method: 无
%% score: 3
%% degree: 750
%% duplicateId: 0
%% body:
1905 年，爱因斯坦创立了“相对论”，提出了著名的质能方程 $E=mc^{2}$。下面涉及到对质能方程理解的几种说法中，正确的是\xzanswer{A}
\fourchoices[answer=A]
{若物体的能量增大，则它的质量增大}
{若物体的能量增大，则它的质量减小}
{若核反应过程质量减小，则需吸收能量}
{若核反应过程质量增大，则会放出能量}
%% answer: A
%% memo:
\memoanswer{
本题原卷及材料中未提供详解，待补充。
}

\item
%% number: 25
%% paperName: 2004年江苏春季理综第 25 题 3 分
%% typeId: 4
%% type: 实验题
%% chapter: 功和能
%% point: DIS研究机械能守恒定律
%% method: 无
%% score: 3
%% degree: 850
%% duplicateId: 0
%% body:
关于“验证机械能守恒定律”的实验，下列叙述中正确的是\xzanswer{C}
\fourchoices[answer=C]
{必须用天平测出重物的质量}
{必须采用直流低压电源}
{必须先接通电源，后释放纸带}
{必须用停表测出重物下落的时间}
%% answer: C
\jdanswer{
C
}
%% memo:
\memoanswer{
本题原卷及材料中未提供详解，待补充。
}

\item
%% number: 32
%% paperName: 2004年江苏春季理综第 32 题 4 分
%% typeId: 3
%% type: 填空题
%% chapter: 机械振动
%% point: 单摆
%% method: 无
%% score: 4
%% degree: 750
%% duplicateId: 0
%% body:
两个单摆甲和乙，它们的摆长之比为 $4:1$，若它们在同一地点做简谐运动，则它们的周期之比 $T_{\text{甲}}:T_{\text{乙}}=$ \tkanswer[1.6cm]{$2:1$}；在甲摆完成 10 次全振动的时间内，乙摆完成的全振动次数为 \tkanswer[1.6cm]{20}。
%% answer: 
%% memo:
\memoanswer{
本题原卷及材料中未提供详解，待补充。
}

\item
%% number: 33
%% paperName: 2004年江苏春季理综第 33 题 4 分
%% typeId: 3
%% type: 填空题
%% chapter: 光学
%% point: 全反射
%% method: 无
%% score: 4
%% degree: 750
%% duplicateId: 0
%% body:
如图\ref{2004江苏春季理综33}所示，水平直线表示空气与某种均匀介质的分界面，与其垂直的为法线。一束单色光从空气射向该介质时，测得入射角 $i=45^{\circ}$，折射角 $r=30^{\circ}$，则该均匀介质的折射率 $n=$ \tkanswer[1.6cm]{$\sqrt{2}$}；若同种单色光从该均匀介质射向空气，则当入射角 $C=$ \tkanswer[1.6cm]{$45^{\circ}$} 时，恰好没有光线射入空气。
\onepicture[num=33, numstyle=arabic, labelprefix={图}, label=2004江苏春季理综33, align=right, w=5.0cm]{figs/33.png}
%% answer: 
%% memo:
\memoanswer{
本题原卷及材料中未提供详解，待补充。
}

\item
%% number: 34
%% paperName: 2004年江苏春季理综第 34 题 7 分
%% typeId: 6
%% type: 计算题
%% chapter: 功和能
%% point: 动能定理
%% method: 无
%% score: 7
%% degree: 650
%% duplicateId: 0
%% body:
质量 $M=6.0\times10^{3}\Ukg$ 的客机，从静止开始沿平直的跑道滑行，当滑行距离 $s=7.2\times10^{2}\Um$ 时，达到起飞速度 $v=60\Ums$。
\begin{enumerate}
	\item
	起飞时飞机的动能多大？
	\item
	若不计滑行过程中所受的阻力，则飞机受到的牵引力为多大？
	\item
	若滑行过程中受到的平均阻力大小为 $F=3.0\times10^{3}\UN$，牵引力与第（2）问中求得的值相等，则要达到上述起飞速度，飞机的滑行距离应为多大？
\end{enumerate}
%% answer: 
\jdanswer{
\begin{enumerate}
	\item
	$E_{k}=1.08\times10^{7}\UJ$

	\item
	$F_{1}=1.5\times10^{4}\UN$

	\item
	$s'=9.0\times10^{2}\Um$

\end{enumerate}
}
%% memo:
\memoanswer{
（1）飞机起飞时的动能为 $E_{k}=\frac{1}{2}Mv^{2}$，代入数值，得

$E_{k}=1.08\times10^{7}\UJ$。

（2）设牵引力为 $F_{1}$，由动能定理，得

$F_{1}s=E_{k}-0$，

代入数值，解得

$F_{1}=1.5\times10^{4}\UN$。

（3）设滑行距离为 $s'$，由动能定理，得

$(F_{1}-F)s'=E_{k}-0$，

整理得

$s'=\frac{E_{k}}{F_{1}-F}$，

代入数值，得

$s'=9.0\times10^{2}\Um$。

评分标准：本题 7 分，其中第（1）、（2）问各 2 分，第（3）问 3 分。其他解法，只要正确可参照以上标准给分。
}

\item
%% number: 35
%% paperName: 2004年江苏春季理综第 35 题 9 分
%% typeId: 6
%% type: 计算题
%% chapter: 电磁感应
%% point: 电磁感应电路分析
%% method: 无
%% score: 9
%% degree: 650
%% duplicateId: 0
%% body:
如图\ref{2004江苏春季理综35}所示，U 形导线框 MNQP 水平放置在磁感应强度 $B=0.2\UT$ 的匀强磁场中，磁感线方向与导线框所在平面垂直，导线 MN 和 PQ 足够长，间距为 $0.5\Um$，横跨在导线框上的导体棒 ab 的电阻 $r=1.0\UO$，接在 NQ 间的电阻 $R=4.0\UO$，电压表为理想电表，其余电阻不计。若导体棒在水平外力作用下以速度 $v=2.0\Ums$ 向左做匀速直线运动，不计导体棒与导线框间的摩擦。
\begin{enumerate}
	\item
	通过电阻 $R$ 的电流方向如何？
	\item
	电压表的示数为多少？
	\item
	若某一时刻撤去水平外力，则从该时刻起，在导体棒运动 $1.0\Um$ 的过程中，通过导体棒的电荷量为多少？
\end{enumerate}
\onepicture[num=35, numstyle=arabic, labelprefix={图}, label=2004江苏春季理综35, align=right, w=7.0cm]{figs/35.png}
%% answer: 
\jdanswer{
\begin{enumerate}
	\item
	$N\to Q$

	\item
	$U=0.16\UV$

	\item
	$Q=2.0\times10^{-2}\UC$

\end{enumerate}
}
%% memo:
\memoanswer{
（1）由右手定则可判断，导体棒中的电流方向为 $b\to a$，则通过电阻 $R$ 的电流方向为 $N\to Q$。

（2）由感应电动势的公式，得

$E=Blv$ ①。

设电路中的电流为 $I$，由闭合电路欧姆定律，得

$I=\frac{E}{R+r}$ ②。

又电压表的示数等于电阻 $R$ 两端的电压值，则有

$U=IR$ ③。

综合①②③式，得

$U=\frac{BlvR}{R+r}$ ④。

代入数值，得

$U=0.16\UV$ ⑤。

（3）撤去水平外力后，导体棒将在安培力的作用下，做减速运动。设在导体棒运动 $x=1.0\Um$ 的过程中，导体棒中产生的感应电动势的平均值为 $E'$，

由法拉第电磁感应定律，得

$E'=\frac{Blx}{\Delta t}$ ⑥。

由闭合电路欧姆定律，得

$I=\frac{E'}{R+r}$ ⑦。

设通过导体棒的电荷量为 $Q$，则有

$Q=I\Delta t$ ⑧。

综合⑥、⑦、⑧式，得

$Q=\frac{Blx}{R+r}$ ⑨。

代入数值，得

$Q=2.0\times10^{-2}\UC$ ⑩。

评分标准：本题 9 分，其中第（1）问 2 分，第（2）问 3 分，第（3）问 4 分。

第（1）问中，正确判断出电阻及中电流方向给 2 分，如果只判断出通过导体棒中电流方向的给 1 分。

第（2）问中，得出④式给 2 分（得出①②③式中的任何一式可给 1 分），得出⑤式再给 1 分。

第（3）问中，得出⑥⑦⑧⑩式各给 1 分。（若将导体棒的运动作匀速运动或匀变速运动处理，则不给分：只写出结果的，也不给分。）

其他解法，只要正确可参照以上标准给分。
}

\item
%% number: 36
%% paperName: 2004年江苏春季理综第 36 题 4 分
%% typeId: 3
%% type: 填空题
%% chapter: 运动和力
%% point: 变加速直线运动
%% method: 无
%% score: 4
%% degree: 750
%% duplicateId: 0
%% body:
设在某次人工降雨中，有一质量恒为 $m$ 的雨滴从高空由静止开始竖直下落，雨滴下落过程中受到的空气阻力大小与下落速度大小成正比，即 $F=kv$，其中 $k$ 为比例系数，则雨滴在刚开始下落的一小段时间内做加速度 \tkanswer[1.6cm]{减小}、速度 \tkanswer[1.6cm]{增大} 的直线运动（以上两空选填“增大”、“减小”或“不变”）。雨滴最终做匀速直线运动的速度表达式为 $v_{m}=$ \tkanswer[1.6cm]{$mg/k$}。
%% answer: 
%% memo:
\memoanswer{
本题原卷及材料中未提供详解，待补充。
}

\end{enumerate}

\end{document}
